
\section*{4. Calcul du @@M0@@ d'un fibré de drapeaux}

4.1.
Soient @@M1@@ un schéma, et @@M2@@ un @@M3@@-Module localement libre de rang @@M4@@. Soit @@M5@@ une suite d'entiers @@M6@@, telle que~:
@@M7@@
Rappelons qu'on appelle foncteur des drapeaux de type @@M8@@ de @@M9@@ le foncteur sur la catégorie des @@M10@@-schémas, à valeurs dans la catégorie des ensembles, qui à un @@M11@@-schéma @@M12@@ associe l'ensemble des suites de quotients de @@M13@@~:
@@M14@@
où, pour tout @@M15@@, @@M16@@ est un quotient de @@M17@@, localement libre de rang @@M18@@.

Un tel foncteur est représentable par un @@M19@@-schéma @@M20@@, qui est muni d'une suite de quotients canonique de l'image inverse de @@M21@@, localement libres de rang @@M22@@

Si @@M23@@, alors @@M24@@.

Si @@M25@@, alors @@M26@@, grassmannienne de type @@M27@@ de @@M28@@.

Si tous les @@M29@@ sont égaux à @@M30@@, on trouve le fibré en drapeau usuel de @@M31@@, noté @@M32@@.

\textbf{Lemme 4.2.} \uline{Soient} @@M33@@ \uline{un schéma}, @@M34@@ \uline{un} @@M35@@\uline{-Module localement libre de rang} @@M36@@, @@M37@@, \uline{avec} @@M38@@ ; \uline{soient} @@M39@@ \uline{deux entiers} @@M40@@, \uline{tels que} @@M41@@ ; \uline{posons} @@M42@@, @@M43@@. \uline{Soient} @@M44@@ \uline{les quotients canoniques de l'image inverse de} @@M45@@ \uline{sur} @@M46@@, \uline{et soit} @@M47@@ \uline{le noyau de} @@M48@@. \uline{Alors on a un} @@M49@@-\uline{isomorphisme canonique}~:
